\documentclass[11pt]{article}
\usepackage[margin=1in]{geometry}
\usepackage{amsmath,amssymb,amsthm,mathtools,microtype,graphicx}
\usepackage[hidelinks]{hyperref}
\newtheorem{theorem}{Theorem}
\newtheorem{lemma}{Lemma}
\newcommand{\R}{\mathbb R}
\newcommand{\C}{\mathbb C}
\newcommand{\V}{\mathcal V}
\newcommand{\dist}{\operatorname{dist}}
\setlength{\parindent}{0pt}
\setlength{\parskip}{0.62em}

\title{Positive $R(\mathcal V)$ Does Not Force Positive $d_{\mathcal V}$\\
\large Counterexample to Remark 2.23 on Fr\'echet approximation schemes}
\author{Run \texttt{fa\_banach\_001}, lane 18}
\date{11 August 2026}

\begin{document}
\maketitle

\begin{center}
\fbox{\parbox{0.91\textwidth}{\textbf{Status.}
Candidate full counterexample, likely valid, for human review. The proof is
self-contained and analytic; novelty remains subject to broader literature
review.}}
\end{center}

\section{Source question}

Let $(X,\|\cdot\|)$ be an infinite-dimensional Fr\'echet space and let
$\V=(V_n)_{n\geq1}$ be nested linear subspaces such that
$\overline{V_n}\subsetneq V_{n+1}$ and
$X=\overline{\bigcup_{n\geq1}V_n}$. Define
\[
 \rho_n(x)=\dist(x,V_n),\qquad
 d_{n,\V}=\sup_{v\in V_{n+1}}\rho_n(v),\qquad
 d_{\V}=\inf_{n\geq1}d_{n,\V},
\]
and
\[
 R(\V)=\inf_{v\in(\bigcup_nV_n)\setminus\{0\}}
          \sup_{t\in\R_+}\|tv\|.
\]
Aksoy and Lewicki prove $d_{\V}>0\Rightarrow R(\V)>0$. They note that the
converse holds when every $V_n$ is finite-dimensional and ask whether the
equivalence holds for arbitrary $\V$ \cite{AL}. The answer is no.

\begin{center}
 \includegraphics[width=0.96\textwidth]{figures/open_problem_crop.png}

 \small Lemma 2.22 and Remark 2.23, source PDF page 16.
\end{center}

\section{Idea of the proof}

Use the usual smooth-function space, but metrize it by a bounded product
F-norm. Along every nonzero ray, every coordinate of this metric saturates,
so the supremum of the norm is exactly one; this forces $R(\V)=1$. For the
subspaces, filter functions by how many derivatives of their germ at zero may
be nonzero. Passing from $V_n$ to $V_{n+1}$ adds just the $n$th jet. A degree
$n$ monomial times a shrinking cutoff removes that jet. Its low derivatives
vanish in the limit, while all uncontrolled high derivatives live in a metric
tail of total weight $2^{-n}$. Thus $d_{n,\V}\leq2^{-n}$.

\section{The counterexample}

Let $X=C^\infty([0,1];\mathbb K)$, where $\mathbb K=\R$ or $\C$, and set
\[
 p_j(f)=\max_{0\leq k\leq j}\|f^{(k)}\|_\infty,
 \qquad
 \|f\|_F=\sum_{j=0}^{\infty}2^{-(j+1)}
          \frac{p_j(f)}{1+p_j(f)}.
 \tag{1}\label{eq:Fnorm}
\]
For $n\geq1$, define
\[
 V_n=\{f\in X:f^{(k)}(0)=0\text{ for every }k\geq n\}.
 \tag{2}\label{eq:Vn}
\]

\begin{theorem}\label{thm:main}
The F-norm in \eqref{eq:Fnorm} makes $X$ a Fr\'echet space, and the subspaces
in \eqref{eq:Vn} are closed, strictly nested, and infinite-dimensional, with
dense union. Nevertheless,
\[
 R(\V)=1\qquad\text{and}\qquad d_{\V}=0.
\]
In particular, $R(\V)>0$ does not imply $d_{\V}>0$ for arbitrary $\V$.
\end{theorem}

\section{Topology and density}

The functions $s\mapsto s/(1+s)$ are increasing and subadditive on
$[0,\infty)$, so \eqref{eq:Fnorm} is an F-norm. It induces the usual
$C^\infty$ topology: convergence in the F-norm forces convergence in every
$p_j$ by looking at the $j$th summand, while convergence in every $p_j$
implies convergence of the sum by splitting it into a finite head and a
uniformly small geometric tail. Hence completeness is the usual completeness
of $C^\infty[0,1]$.

Every jet functional $f\mapsto f^{(k)}(0)$ is continuous. Thus $V_n$, an
intersection of their kernels, is closed. We have $V_n\subset V_{n+1}$, and
the inclusion is proper because $x\mapsto x^n$ lies in $V_{n+1}\setminus
V_n$. Each $V_n$ is infinite-dimensional because it contains every smooth
function supported in a compact subset of $(0,1]$.

The union contains every polynomial: a polynomial of degree $d$ belongs to
$V_{d+1}$. For completeness, we verify that polynomials are dense in the full
smooth topology. Fix $f\in X$, $m\geq1$, and $\varepsilon>0$. By the
Weierstrass theorem choose a polynomial $q$ with
$\|f^{(m)}-q\|_\infty<\varepsilon$. Define the polynomial
\[
 P(x)=\sum_{r=0}^{m-1}\frac{f^{(r)}(0)}{r!}x^r
 +\int_0^x\frac{(x-s)^{m-1}}{(m-1)!}\,q(s)\,ds.
\]
Then $P^{(m)}=q$ and $P^{(r)}(0)=f^{(r)}(0)$ for $r<m$. Repeated integration
gives, for $0\leq r<m$,
\[
 \|f^{(r)}-P^{(r)}\|_\infty
 \leq\frac{\|f^{(m)}-q\|_\infty}{(m-r)!}<\varepsilon.
\]
Thus $p_m(f-P)<\varepsilon$. The case $m=0$ is Weierstrass approximation
itself. Hence $X=\overline{\bigcup_nV_n}$.

\section{Exact value of \texorpdfstring{$R(\mathcal V)$}{R(V)}}

If $f\neq0$, then $p_j(f)\geq p_0(f)>0$ for all $j$. For $t\geq0$,
\[
 \|tf\|_F=\sum_{j=0}^{\infty}2^{-(j+1)}
 \frac{t p_j(f)}{1+t p_j(f)}\leq1.
\]
Every summand increases to its full weight as $t\to\infty$. Monotone
convergence therefore yields
\[
 \lim_{t\to\infty}\|tf\|_F
 =\sum_{j=0}^{\infty}2^{-(j+1)}=1.
\]
Consequently $\sup_{t\in\R_+}\|tf\|_F=1$ for every nonzero $f$, in
particular for every nonzero $f\in\bigcup_nV_n$. Hence
\[
 R(\V)=1. \tag{3}\label{eq:R}
\]

\section{Decay of \texorpdfstring{$d_{n,\mathcal V}$}{d(n,V)}}

Fix $n\geq1$ and $f\in V_{n+1}$. Put $a=f^{(n)}(0)$. Choose
$\chi\in C^\infty([0,\infty))$ such that $\chi=1$ near zero and
$\chi=0$ on $[1,\infty)$. For $0<\delta\leq1$, define
\[
 h_\delta(x)=a\frac{x^n}{n!}\chi(x/\delta).
 \tag{4}\label{eq:cutoff}
\]
Near zero, $h_\delta(x)=a x^n/n!$. Therefore
\[
 h_\delta^{(n)}(0)=a,qquad
 h_\delta^{(k)}(0)=0\quad(k>n).
\]
Because $f\in V_{n+1}$ already has $f^{(k)}(0)=0$ for $k>n$, it follows that
$g_\delta=f-h_\delta\in V_n$.

For every fixed $k<n$, Leibniz' rule and the support condition give a constant
$C_{n,k,\chi}$, independent of $a$ and $\delta$, such that
\[
 \|h_\delta^{(k)}\|_\infty
 \leq C_{n,k,\chi}|a|\delta^{,n-k}.
 \tag{5}\label{eq:lowder}
\]
Indeed, each term is a constant multiple of
$x^{n-r}\delta^{-(k-r)}\chi^{(k-r)}(x/\delta)$ on $0\leq x\leq\delta$,
and hence has size $O(\delta^{n-k})$. It follows that
$p_j(h_\delta)\to0$ for every $j<n$.

Using $s/(1+s)\leq1$ on the remaining terms, we obtain
\begin{align*}
 \rho_n(f)
 &\leq \|f-g_\delta\|_F=\|h_\delta\|_F,\\
 \limsup_{\delta\downarrow0}\|h_\delta\|_F
 &\leq \sum_{j=n}^{\infty}2^{-(j+1)}=2^{-n}.
\end{align*}
Thus $\rho_n(f)\leq2^{-n}$ for every $f\in V_{n+1}$. Taking the supremum,
\[
 d_{n,\V}\leq2^{-n};
\]
and therefore
\[
 d_{\V}=\inf_{n\geq1}d_{n,\V}=0. \tag{6}\label{eq:d}
\]
Equations \eqref{eq:R} and \eqref{eq:d} prove Theorem~\ref{thm:main}.

\section{Verification, scope, and novelty status}

The proof uses no computation. The only limiting estimate relevant to the
supremum in $d_{n,\V}$ is uniform in the required direction: for every fixed
$f\in V_{n+1}$ it proves the same upper bound $2^{-n}$, even though the
cutoff scale may depend on $f^{(n)}(0)$. If that jet is zero, then
$f\in V_n$ already.

The construction works over both real and complex scalars. All $V_n$ are
closed and infinite-dimensional, so the failure is precisely outside the
finite-dimensional hypothesis highlighted in Remark 2.23. It does not
contradict the proved implication $d_{\V}>0\Rightarrow R(\V)>0$.

A bounded search of the run indexes and exact web queries found the source
paper and its journal version, but no later resolution of Remark 2.23.
Priority is not asserted. Expert proof review and a broader citation search
are recommended before dissemination.

\begin{thebibliography}{9}
\bibitem{AL}
A.~G.~Aksoy and G.~Lewicki,
\emph{Bernstein's Lethargy Theorem in Fr\'echet Spaces},
arXiv:1503.06190 (2015).
\end{thebibliography}

\end{document}
